\begin{figure*}[tbp]
\centering
\begin{minipage}[t]{0.72\textwidth}\vspace{0pt}
\centering
\begin{tikzpicture}[baseline=(current bounding box.north)]
\begin{axis}[
  width=\linewidth, height=0.25\textheight,
  date coordinates in=x, xticklabel=\day.\month, xticklabel style={rotate=45,anchor=east},
  xlabel={Date/Time (UTC)}, ylabel={SST ($^\circ$C)},
  axis lines=left,
  grid=major, grid style={draw=CUDAHeatGrid},
  xtick distance=1, minor x tick num=1,
  ymin=26.93553, ymax=29.797368,
  legend style={at={(0.5,1.03)},anchor=south,legend columns=3,draw=none,font=\small},
]
\addplot[SSTLine,thick,no markers] table[x=time,y=sst,col sep=comma] {tikz/events/event-000002_point-0.csv};
\addlegendentry{SST}
\addplot[ThresholdLine,thin,opacity=1,no markers] table[x=time,y=threshold,col sep=comma] {tikz/events/event-000002_point-0.csv};
\addlegendentry{Threshold}
\addplot[ClimatologyLine,thin,no markers] table[x=time,y=climatology,col sep=comma] {tikz/events/event-000002_point-0.csv};
\addlegendentry{Climatology}
\end{axis}
\end{tikzpicture}
\end{minipage}\hfill
\begin{minipage}[t]{0.24\textwidth}\vspace{0pt}
\footnotesize
\textbf{Detection settings}\par\vspace{2pt}
\begin{tabular}{@{}ll@{}}
\textcolor{CUDAHeatSlate}{Duration} & 142 h \\
\textcolor{CUDAHeatSlate}{Baseline} & 2019–2023 \\
\textcolor{CUDAHeatSlate}{Percentile} & 90\% \\
\textcolor{CUDAHeatSlate}{Internal gap} & 0 h \\
\textcolor{CUDAHeatSlate}{Join gap} & 12 h \\
\end{tabular}
\par\vspace{5pt}\textbf{Normalized metrics}\par\vspace{2pt}
\begin{tikzpicture}
\begin{axis}[
  xbar, width=\linewidth, height=3.10cm,
  xmin=0, xmax=1, xtick={0,0.5,1}, minor x tick num=1,
  ytick={0,1,2,3,4}, yticklabels={},
  y dir=reverse, enlarge y limits=0.30,
  axis x line*=bottom, axis y line=none,
  ytick style={draw=none}, clip=false,
  tick label style={font=\scriptsize},
  xmajorgrids, grid style={draw=CUDAHeatGrid},
  bar width=3.5pt,
]
\addplot[draw=none,fill=gray!60] coordinates {(1.0000,0) (0.2651,1) (0.0000,2) (0.0000,3) (0.0000,4)};
\node[anchor=east,xshift=-3pt,font=\scriptsize,text=CUDAHeatSlate] at (axis cs:0,0) {$F_T$};
\node[anchor=east,xshift=-3pt,font=\scriptsize,text=CUDAHeatSlate] at (axis cs:0,1) {$DSEF$};
\node[anchor=east,xshift=-3pt,font=\scriptsize,text=CUDAHeatSlate] at (axis cs:0,2) {$NBCF$};
\node[anchor=east,xshift=-3pt,font=\scriptsize,text=CUDAHeatSlate] at (axis cs:0,3) {$NBTF$};
\node[anchor=east,xshift=-3pt,font=\scriptsize,text=CUDAHeatSlate] at (axis cs:0,4) {$NDHR$};
\end{axis}
\end{tikzpicture}
\par\vspace{5pt}\textbf{Dimensional metrics}\par\vspace{2pt}
\begin{tabular}{@{}ll@{}}
\textcolor{CUDAHeatSlate}{$HC$} & 56.313 \\
\end{tabular}
\end{minipage}
\caption{Hourly SST, climatological threshold and climatological mean for CUDAHeat event 2 (run \texttt{7b6805fd-6c65-406b-8916-7237c2a0b469}). Event period: August 2022 to September 2022.}
\label{fig:event:7b6805fd-6c65-406b-8916-7237c2a0b469:2:point-0}
\end{figure*}
